\begin{table}[t]
\centering
\caption{How much larger is the HEALPix tile, once compressed?}
\label{tab:headline}
\scriptsize
\setlength{\tabcolsep}{4pt}
\begin{tabular}{@{}lrlrrr@{}}
\toprule
Target & Error & \multicolumn{3}{c}{HEALPix / UTM size, by transfer} & \\
\cmidrule(lr){3-5}
RMSE & (\% contrast) & C--T [class range] & PSF & IDW & $n$ \\
\midrule
uncompressed & -- & 2.58 & 2.58 & 2.58 & -- \\
\midrule
$10^{-4}$ & 0.25 & 2.27 [2.18--2.33] & 2.62 & 2.49 & 39 \\
$3\times10^{-4}$ & 0.75 & 2.18 [1.96--2.23] & 2.64 & 2.51 & 39 \\
$10^{-3}$ & 2.5 & 1.82 [1.58--2.03] & 2.65 & 2.53 & 39 \\
$3\times10^{-3}$ & 7.5 & 1.61 [1.37--1.90] & 2.60 & 2.47 & 38 \\
$10^{-2}$ & 25 & 1.53 [1.28--1.60] & 2.32 & 2.07 & 37 \\
\bottomrule
\end{tabular}
\vspace{3pt}
\begin{minipage}{\columnwidth}
\footnotesize
Median over the 40 regions of the compressed size of the HEALPix tile divided by that of the native UTM crop of the same nominal ground area, at the same reconstruction error and with no storage overhead charged to either side ($\gamma=1$); a value above one means the HEALPix tile is larger. The second column expresses each target as a fraction of a typical scene's contrast (the median standard deviation of the 40 tiles, 0.040 in reflectance), so that $10^{-2}$ is an error of about a quarter of the contrast and $10^{-4}$ a quarter of a percent. The C--T (Clough--Tocher) column is the carrier the paper uses, with the range of the four land-cover class medians in brackets; the PSF (PSF-aware) and IDW (inverse-distance) columns are the same tile rebuilt with the other two transfers of Section~\ref{sec:transfer_cost}, on the $n$ regions all three sweeps bracket. Under an operational overhead $\gamma$ on the native product every entry divides by $\gamma$ (1.3 in the illustrative scenario of Section~\ref{sec:rd_illustrative}).
\end{minipage}
\end{table}
